
\documentclass{article}
\usepackage{colortbl}
\usepackage{makecell}
\usepackage{multirow}
\usepackage{supertabular}

\begin{document}

\newcounter{utterance}

\centering \large Interaction Transcript for game `clean\_up', experiment `2\_hard\_3obj\_nl', episode 0 with deepseek{-}v4{-}pro.
\vspace{24pt}

{ \footnotesize  \setcounter{utterance}{1}
\setlength{\tabcolsep}{0pt}
\begin{supertabular}{c@{$\;$}|p{.15\linewidth}@{}p{.15\linewidth}p{.15\linewidth}p{.15\linewidth}p{.15\linewidth}p{.15\linewidth}}
    \# & $\;$A & \multicolumn{4}{c}{Game Master} & $\;\:$B\\
    \hline

    \theutterance \stepcounter{utterance}  
    & & \multicolumn{4}{p{0.6\linewidth}}{
        \cellcolor[rgb]{0.9,0.9,0.9}{
            \makecell[{{p{\linewidth}}}]{
                \texttt{\tiny{[P1$\langle$GM]}}
                \texttt{Ik ben uw spelleider, en u speelt een samenwerkingsspel met het volgende raster als speelbord:} \\
\\ 
\texttt{```} \\
\texttt{ 1234567} \\
\texttt{╔══╤════╗} \\
\texttt{║◌◌│C◌◌◌║ 1} \\
\texttt{╟──┤◌L┌─╢ 2} \\
\texttt{║◌◌├──┤◌║ 3} \\
\texttt{╟──┘◌◌├─╢ 4} \\
\texttt{║◌◌P◌◌│◌║ 5} \\
\texttt{╟─┐◌┌─┼─╢ 6} \\
\texttt{║◌│◌│◌│◌║ 7} \\
\texttt{╚═╧═╧═╧═╝} \\
\\ 
\texttt{```} \\
\\ 
\texttt{* De bovenrand toont x{-}coördinaten die naar rechts oplopen, en de rechterrand y{-}coördinaten die naar beneden oplopen.} \\
\texttt{* De volgende objecten zijn willekeurig op uw raster geplaatst: 'C', 'L', 'P'.} \\
\\ 
\\ 
\texttt{De andere speler ziet een variatie van het speelbord, waarbij de objecten op verschillende willekeurige locaties zijn geplaatst. U kunt het bord van de andere speler niet zien en zij kunnen het uwe niet zien.} \\
\\ 
\\ 
\texttt{**Doel:**} \\
\\ 
\texttt{Beide spelers moeten de objecten op hun respectievelijke achtergrond zo verplaatsen dat identieke objecten op dezelfde coördinaten terechtkomen. U moet communiceren met de andere speler om een gemeenschappelijke doelconfiguratie overeen te komen.} \\
\\ 
\\ 
\texttt{**Regels:**} \\
\\ 
\texttt{* In elke beurt kunt u precies één van de volgende twee commando's sturen:} \\
\texttt{ 1. `ZEG: <BERICHT>`: om een bericht (alles tot de volgende regeleinde) naar de andere speler te sturen. Ik zal het doorsturen naar uw partner.} \\
\texttt{2. `BEWEEG: <OBJECT>, (<X>, <Y>)`: om een object naar een nieuwe positie te verplaatsen, waarbij `<X>` de kolom is en `<Y>` de rij. Ik zal u laten weten of uw zet succesvol was of niet.} \\
\texttt{* Als u zich niet aan het formaat houdt, of meerdere commando's tegelijk verstuurt, moet ik u straffen.} \\
\texttt{* Als beide spelers meer dan 8 strafpunten verzamelen, verliezen jullie allebei het spel.} \\
\texttt{* Het is van vitaal belang dat u met de andere speler communiceert over uw doelstatus! De *enige* manier waarop u uw strategie aan de andere speler kunt overbrengen is door gebruik te maken van het `ZEG: <BERICHT>` commando!} \\
\\ 
\\ 
\texttt{**Objecten verplaatsen**} \\
\\ 
\texttt{* U kunt alleen objecten verplaatsen naar cellen binnen de grenzen van het raster. De doelcel moet leeg zijn, d.w.z. deze mag alleen het symbool '◌' bevatten.} \\
\texttt{* Als u probeert een object te verplaatsen naar een plek die niet leeg is, of als u probeert het buiten het raster te verplaatsen, moet ik u straffen. U krijgt een nieuwe poging.} \\
\texttt{* Voordat u een zet doet, moet u controleren of de doelplek leeg is, en geen letter, kader of lijn bevat!} \\
\\ 
\\ 
\texttt{**Einde van het spel**} \\
\\ 
\texttt{Als u denkt dat u het doel bereikt hebt om alle objecten uit te lijnen, kunt u de andere speler vragen het spel te beëindigen door `ZEG: klaar?` te sturen. Als de andere speler u vraagt om het spel te beëindigen en u antwoordt met `ZEG: klaar!`, zal het spel eindigen.} \\
\\ 
\texttt{Beide spelers winnen als het spel eindigt binnen 12 rondes, waarbij één ronde gedefinieerd wordt als twee spelers die elk een geldig commando sturen.} \\
\\ 
\texttt{**Score:**} \\
\\ 
\texttt{Hoe dichter de identieke objecten bij elkaar liggen op beide speelborden, hoe meer punten u krijgt. Strafpunten verminderen uw punten. Kunt u het record verslaan?} \\
\\ 
\\ 
\texttt{Stuur een bericht naar de andere speler om het spel te starten!} \\
            }
        }
    }
    & \\ \\

    \theutterance \stepcounter{utterance}  
    & \multicolumn{4}{p{0.6\linewidth}}{
        \cellcolor[rgb]{0.9,0.9,0.9}{
            \makecell[{{p{\linewidth}}}]{
                \texttt{\tiny{[P1$\rangle$GM]}}
                \texttt{ZEG: Hallo! Laten we onze objectposities delen. Mijn bord heeft:} \\
\texttt{C op (3,1)} \\
\texttt{L op (3,2)} \\
\texttt{P op (3,5)} \\
\\ 
\texttt{Wat zijn de posities van C, L en P op jouw bord?} \\
            }
        }
    }
    & & \\ \\

    \theutterance \stepcounter{utterance}  
    & & & \multicolumn{2}{p{0.3\linewidth}}{
        \cellcolor[rgb]{0.9,0.9,0.9}{
            \makecell[{{p{\linewidth}}}]{
                \texttt{\tiny{[GM$|$GM]}}
                \texttt{ZEG: Hallo! Laten we onze objectposities delen. Mijn bord heeft:} \\
\texttt{C op (3,1)} \\
\texttt{L op (3,2)} \\
\texttt{P op (3,5)} \\
\\ 
\texttt{Wat zijn de posities van C, L en P op jouw bord?} \\
            }
        }
    }
    & & \\ \\

    \theutterance \stepcounter{utterance}  
    & & & \multicolumn{2}{p{0.3\linewidth}}{
        \cellcolor[rgb]{0.9,0.9,0.9}{
            \makecell[{{p{\linewidth}}}]{
                \texttt{\tiny{[GM$|$GM]}}
                \texttt{Invalid format: tail is not empty: '} \\
\texttt{C op (3,1)} \\
\texttt{L op (3,2)} \\
\texttt{P op (3,5)} \\
\\ 
\texttt{Wat zijn de posities van C, L en P op jouw bord?} \\
            }
        }
    }
    & & \\ \\

    \theutterance \stepcounter{utterance}  
    & & \multicolumn{4}{p{0.6\linewidth}}{
        \cellcolor[rgb]{0.9,0.9,0.9}{
            \makecell[{{p{\linewidth}}}]{
                \texttt{\tiny{[P1$\langle$GM]}}
                \texttt{Straf: Je bericht mag niets na het commando bevatten!} \\
\\ 
\texttt{Zorg ervoor dat je antwoord alleen `ZEG: <BERICHT>` of `BEWEEG: <OBJECT>, (<X>, <Y>)` bevat, en niets anders!} \\
\\ 
\texttt{Jullie hebben gezamenlijk 1 van 8 straffen verzameld.} \\
\texttt{Probeer het opnieuw!} \\
            }
        }
    }
    & \\ \\

    \theutterance \stepcounter{utterance}  
    & \multicolumn{4}{p{0.6\linewidth}}{
        \cellcolor[rgb]{0.9,0.9,0.9}{
            \makecell[{{p{\linewidth}}}]{
                \texttt{\tiny{[P1$\rangle$GM]}}
                \texttt{ZEG: Hallo! Laten we onze posities delen. Mijn C staat op (3,1), L op (3,2), P op (3,5). Waar staan die van jou?} \\
            }
        }
    }
    & & \\ \\

    \theutterance \stepcounter{utterance}  
    & & & \multicolumn{2}{p{0.3\linewidth}}{
        \cellcolor[rgb]{0.9,0.9,0.9}{
            \makecell[{{p{\linewidth}}}]{
                \texttt{\tiny{[GM$|$GM]}}
                \texttt{ZEG: Hallo! Laten we onze posities delen. Mijn C staat op (3,1), L op (3,2), P op (3,5). Waar staan die van jou?} \\
            }
        }
    }
    & & \\ \\

    \theutterance \stepcounter{utterance}  
    & & \multicolumn{4}{p{0.6\linewidth}}{
        \cellcolor[rgb]{0.9,0.9,0.9}{
            \makecell[{{p{\linewidth}}}]{
                \texttt{\tiny{[GM$\rangle$P2]}}
                \texttt{Ik ben uw spelleider, en u speelt een samenwerkingsspel met het volgende raster als speelbord:} \\
\\ 
\texttt{```} \\
\texttt{ 1234567} \\
\texttt{╔══╤════╗} \\
\texttt{║◌◌│L◌◌C║ 1} \\
\texttt{╟──┤◌◌┌─╢ 2} \\
\texttt{║◌◌├──┤◌║ 3} \\
\texttt{╟──┘◌◌├─╢ 4} \\
\texttt{║◌◌P◌◌│◌║ 5} \\
\texttt{╟─┐◌┌─┼─╢ 6} \\
\texttt{║◌│◌│◌│◌║ 7} \\
\texttt{╚═╧═╧═╧═╝} \\
\\ 
\texttt{```} \\
\\ 
\texttt{* De bovenrand toont x{-}coördinaten die naar rechts oplopen, en de rechterrand y{-}coördinaten die naar beneden oplopen.} \\
\texttt{* De volgende objecten zijn willekeurig op uw raster geplaatst: 'C', 'L', 'P'.} \\
\\ 
\\ 
\texttt{De andere speler ziet een variatie van het speelbord, waarbij de objecten op verschillende willekeurige locaties zijn geplaatst. U kunt het bord van de andere speler niet zien en zij kunnen het uwe niet zien.} \\
\\ 
\\ 
\texttt{**Doel:**} \\
\\ 
\texttt{Beide spelers moeten de objecten op hun respectievelijke achtergrond zo verplaatsen dat identieke objecten op dezelfde coördinaten terechtkomen. U moet communiceren met de andere speler om een gemeenschappelijke doelconfiguratie overeen te komen.} \\
\\ 
\\ 
\texttt{**Regels:**} \\
\\ 
\texttt{* In elke beurt kunt u precies één van de volgende twee commando's sturen:} \\
\texttt{ 1. `ZEG: <BERICHT>`: om een bericht (alles tot de volgende regeleinde) naar de andere speler te sturen. Ik zal het doorsturen naar uw partner.} \\
\texttt{2. `BEWEEG: <OBJECT>, (<X>, <Y>)`: om een object naar een nieuwe positie te verplaatsen, waarbij `<X>` de kolom is en `<Y>` de rij. Ik zal u laten weten of uw zet succesvol was of niet.} \\
\texttt{* Als u zich niet aan het formaat houdt, of meerdere commando's tegelijk verstuurt, moet ik u straffen.} \\
\texttt{* Als beide spelers meer dan 8 strafpunten verzamelen, verliezen jullie allebei het spel.} \\
\texttt{* Het is van vitaal belang dat u met de andere speler communiceert over uw doelstatus! De *enige* manier waarop u uw strategie aan de andere speler kunt overbrengen is door gebruik te maken van het `ZEG: <BERICHT>` commando!} \\
\\ 
\\ 
\texttt{**Objecten verplaatsen**} \\
\\ 
\texttt{* U kunt alleen objecten verplaatsen naar cellen binnen de grenzen van het raster. De doelcel moet leeg zijn, d.w.z. deze mag alleen het symbool '◌' bevatten.} \\
\texttt{* Als u probeert een object te verplaatsen naar een plek die niet leeg is, of als u probeert het buiten het raster te verplaatsen, moet ik u straffen. U krijgt een nieuwe poging.} \\
\texttt{* Voordat u een zet doet, moet u controleren of de doelplek leeg is, en geen letter, kader of lijn bevat!} \\
\\ 
\\ 
\texttt{**Einde van het spel**} \\
\\ 
\texttt{Als u denkt dat u het doel bereikt hebt om alle objecten uit te lijnen, kunt u de andere speler vragen het spel te beëindigen door `ZEG: klaar?` te sturen. Als de andere speler u vraagt om het spel te beëindigen en u antwoordt met `ZEG: klaar!`, zal het spel eindigen.} \\
\\ 
\texttt{Beide spelers winnen als het spel eindigt binnen 12 rondes, waarbij één ronde gedefinieerd wordt als twee spelers die elk een geldig commando sturen.} \\
\\ 
\texttt{**Score:**} \\
\\ 
\texttt{Hoe dichter de identieke objecten bij elkaar liggen op beide speelborden, hoe meer punten u krijgt. Strafpunten verminderen uw punten. Kunt u het record verslaan?} \\
\\ 
\texttt{De andere speler begon het spel door dit bericht te sturen:} \\
\\ 
\texttt{"Hallo! Laten we onze posities delen. Mijn C staat op (3,1), L op (3,2), P op (3,5). Waar staan die van jou?"} \\
\\ 
\texttt{Wat is uw eerste opdracht?} \\
            }
        }
    }
    & \\ \\

    \theutterance \stepcounter{utterance}  
    & & & \multicolumn{4}{p{0.6\linewidth}}{
        \cellcolor[rgb]{0.9,0.9,0.9}{
            \makecell[{{p{\linewidth}}}]{
                \texttt{\tiny{[GM$\langle$P2]}}
                \texttt{Ik zal eerst mijn positie delen met de andere speler.} \\
\\ 
\texttt{Mijn objecten staan op:} \\
\texttt{{-} C op (7,1)} \\
\texttt{{-} L op (2,1)} \\
\texttt{{-} P op (3,5)} \\
\\ 
\texttt{```} \\
\texttt{ZEG: Hallo! Mijn C staat op (7,1), L op (2,1), P op (3,5). Jouw P en mijn P staan al op dezelfde positie! We moeten C en L nog uitlijnen. Laten we afspreken welke posities we gebruiken. Welke posities voor C en L werken voor jou?} \\
\texttt{```} \\
            }
        }
    }
     \\ \\

    \theutterance \stepcounter{utterance}  
    & & & \multicolumn{2}{p{0.3\linewidth}}{
        \cellcolor[rgb]{0.9,0.9,0.9}{
            \makecell[{{p{\linewidth}}}]{
                \texttt{\tiny{[GM$|$GM]}}
                \texttt{Ik zal eerst mijn positie delen met de andere speler.} \\
\\ 
\texttt{Mijn objecten staan op:} \\
\texttt{{-} C op (7,1)} \\
\texttt{{-} L op (2,1)} \\
\texttt{{-} P op (3,5)} \\
\\ 
\texttt{```} \\
\texttt{ZEG: Hallo! Mijn C staat op (7,1), L op (2,1), P op (3,5). Jouw P en mijn P staan al op dezelfde positie! We moeten C en L nog uitlijnen. Laten we afspreken welke posities we gebruiken. Welke posities voor C en L werken voor jou?} \\
\texttt{```} \\
            }
        }
    }
    & & \\ \\

    \theutterance \stepcounter{utterance}  
    & & & \multicolumn{2}{p{0.3\linewidth}}{
        \cellcolor[rgb]{0.9,0.9,0.9}{
            \makecell[{{p{\linewidth}}}]{
                \texttt{\tiny{[GM$|$GM]}}
                \texttt{Invalid format: head is not empty: 'Ik zal eerst mijn positie delen met de andere speler.} \\
\\ 
\texttt{Mijn objecten staan op:} \\
\texttt{{-} C op (7,1)} \\
\texttt{{-} L op (2,1)} \\
\texttt{{-} P op (3,5)} \\
\\ 
\\ 
\texttt{'} \\
            }
        }
    }
    & & \\ \\

    \theutterance \stepcounter{utterance}  
    & & \multicolumn{4}{p{0.6\linewidth}}{
        \cellcolor[rgb]{0.9,0.9,0.9}{
            \makecell[{{p{\linewidth}}}]{
                \texttt{\tiny{[GM$\rangle$P2]}}
                \texttt{Straf: Je bericht mag niets vóór het commando bevatten!} \\
\\ 
\texttt{Zorg ervoor dat je antwoord alleen `ZEG: <BERICHT>` of `BEWEEG: <OBJECT>, (<X>, <Y>)` bevat, en niets anders!} \\
\\ 
\texttt{Jullie hebben gezamenlijk 2 van 8 straffen verzameld.} \\
\texttt{Probeer het opnieuw!} \\
            }
        }
    }
    & \\ \\

    \theutterance \stepcounter{utterance}  
    & & & \multicolumn{4}{p{0.6\linewidth}}{
        \cellcolor[rgb]{0.9,0.9,0.9}{
            \makecell[{{p{\linewidth}}}]{
                \texttt{\tiny{[GM$\langle$P2]}}
                \texttt{ZEG: Hallo! Mijn C staat op (7,1), L op (2,1), P op (3,5). Jouw P en mijn P staan al op dezelfde positie! We moeten C en L nog uitlijnen. Laten we afspreken welke posities we gebruiken. Welke posities voor C en L werken voor jou?} \\
            }
        }
    }
     \\ \\

    \theutterance \stepcounter{utterance}  
    & & & \multicolumn{2}{p{0.3\linewidth}}{
        \cellcolor[rgb]{0.9,0.9,0.9}{
            \makecell[{{p{\linewidth}}}]{
                \texttt{\tiny{[GM$|$GM]}}
                \texttt{ZEG: Hallo! Mijn C staat op (7,1), L op (2,1), P op (3,5). Jouw P en mijn P staan al op dezelfde positie! We moeten C en L nog uitlijnen. Laten we afspreken welke posities we gebruiken. Welke posities voor C en L werken voor jou?} \\
            }
        }
    }
    & & \\ \\

    \theutterance \stepcounter{utterance}  
    & & \multicolumn{4}{p{0.6\linewidth}}{
        \cellcolor[rgb]{0.9,0.9,0.9}{
            \makecell[{{p{\linewidth}}}]{
                \texttt{\tiny{[P1$\langle$GM]}}
                \texttt{Je bericht is doorgestuurd naar de andere speler.} \\
\\ 
\\ 
\texttt{Je speelt momenteel ronde 1 van maximaal 12.} \\
\texttt{Jullie hebben gezamenlijk 2 van 8 straffen verzameld.} \\
\texttt{De andere speler stuurde dit bericht:} \\
\\ 
\\ 
\texttt{"Hallo! Mijn C staat op (7,1), L op (2,1), P op (3,5). Jouw P en mijn P staan al op dezelfde positie! We moeten C en L nog uitlijnen. Laten we afspreken welke posities we gebruiken. Welke posities voor C en L werken voor jou?"} \\
\\ 
\\ 
\texttt{Wat is je volgende opdracht?} \\
            }
        }
    }
    & \\ \\

    \theutterance \stepcounter{utterance}  
    & \multicolumn{4}{p{0.6\linewidth}}{
        \cellcolor[rgb]{0.9,0.9,0.9}{
            \makecell[{{p{\linewidth}}}]{
                \texttt{\tiny{[P1$\rangle$GM]}}
                \texttt{ZEG: We hebben geluk met P op (3,5)! Voor C en L: aangezien jouw L op (2,1) staat en mijn L op (3,2), en mijn C op (3,1) en jouw C op (7,1), stel ik voor dat we C op (2,1) zetten en L op (3,1). Dan moet ik mijn C naar (2,1) verplaatsen en L naar (3,1). Jij zou jouw L op (2,1) kunnen laten staan en jouw C van (7,1) naar (3,1) verplaatsen. Wat vind je daarvan?} \\
            }
        }
    }
    & & \\ \\

    \theutterance \stepcounter{utterance}  
    & & & \multicolumn{2}{p{0.3\linewidth}}{
        \cellcolor[rgb]{0.9,0.9,0.9}{
            \makecell[{{p{\linewidth}}}]{
                \texttt{\tiny{[GM$|$GM]}}
                \texttt{ZEG: We hebben geluk met P op (3,5)! Voor C en L: aangezien jouw L op (2,1) staat en mijn L op (3,2), en mijn C op (3,1) en jouw C op (7,1), stel ik voor dat we C op (2,1) zetten en L op (3,1). Dan moet ik mijn C naar (2,1) verplaatsen en L naar (3,1). Jij zou jouw L op (2,1) kunnen laten staan en jouw C van (7,1) naar (3,1) verplaatsen. Wat vind je daarvan?} \\
            }
        }
    }
    & & \\ \\

    \theutterance \stepcounter{utterance}  
    & & \multicolumn{4}{p{0.6\linewidth}}{
        \cellcolor[rgb]{0.9,0.9,0.9}{
            \makecell[{{p{\linewidth}}}]{
                \texttt{\tiny{[GM$\rangle$P2]}}
                \texttt{Je bericht is doorgestuurd naar de andere speler.} \\
\\ 
\\ 
\texttt{Je speelt momenteel ronde 2 van maximaal 12.} \\
\texttt{Jullie hebben gezamenlijk 2 van 8 straffen verzameld.} \\
\texttt{De andere speler stuurde dit bericht:} \\
\\ 
\\ 
\texttt{"We hebben geluk met P op (3,5)! Voor C en L: aangezien jouw L op (2,1) staat en mijn L op (3,2), en mijn C op (3,1) en jouw C op (7,1), stel ik voor dat we C op (2,1) zetten en L op (3,1). Dan moet ik mijn C naar (2,1) verplaatsen en L naar (3,1). Jij zou jouw L op (2,1) kunnen laten staan en jouw C van (7,1) naar (3,1) verplaatsen. Wat vind je daarvan?"} \\
\\ 
\\ 
\texttt{Wat is je volgende opdracht?} \\
            }
        }
    }
    & \\ \\

    \theutterance \stepcounter{utterance}  
    & & & \multicolumn{4}{p{0.6\linewidth}}{
        \cellcolor[rgb]{0.9,0.9,0.9}{
            \makecell[{{p{\linewidth}}}]{
                \texttt{\tiny{[GM$\langle$P2]}}
                \texttt{Ik analyseer het voorstel van de andere speler:} \\
\\ 
\texttt{{-} Doel: C op (2,1), L op (3,1), P op (3,5)} \\
\texttt{{-} Mijn huidige posities: C op (7,1), L op (2,1), P op (3,5)} \\
\texttt{{-} Mijn L staat al op (2,1), maar volgens het voorstel moet L naar (3,1). Mijn C moet van (7,1) naar (2,1) of (3,1)?} \\
\\ 
\texttt{Wacht, even goed kijken. Het voorstel is C op (2,1) en L op (3,1). Mijn L staat op (2,1) en mijn C op (7,1). Dus ik moet L verplaatsen van (2,1) naar (3,1), en C van (7,1) naar (2,1). Dat klopt.} \\
\\ 
\texttt{Maar ik moet eerst controleren of (2,1) en (3,1) leeg zijn. Mijn L staat op (2,1), dus als ik L verplaats naar (3,1), dan komt (2,1) vrij voor C. (3,1) moet leeg zijn. Laat ik het bord checken: op rij 1 zijn de cellen (1,1) en (2,1) vóór de │. Mijn L staat op (2,1), C op (7,1). Positie (3,1) {-} even kijken naar het bord. De bovenste rij is: ╔══╤════╗ met rij 1: ║◌◌│L◌◌C║. Dus (1,1)=◌, (2,1)=◌ maar daar staat L, (3,1)=│ (muur), (4,1)=L? Nee, (2,1)=L, daarna │, dan (4,1)=◌, (5,1)=◌, (6,1)=C? Even opnieuw.} \\
\\ 
\texttt{Rij 1: ║◌◌│L◌◌C║} \\
\texttt{{-} (1,1) = ◌} \\
\texttt{{-} (2,1) = ◌ maar hier staat L? Nee, "◌◌│L◌◌C" {-} de eerste twee zijn ◌, dan │, dan L, ◌, ◌, C. Dus:} \\
\texttt{{-} (1,1) = ◌} \\
\texttt{{-} (2,1) = ◌} \\
\texttt{{-} (3,1) = │ (muur)} \\
\texttt{{-} (4,1) = L} \\
\texttt{{-} (5,1) = ◌} \\
\texttt{{-} (6,1) = ◌} \\
\texttt{{-} (7,1) = C} \\
\\ 
\texttt{Oh! Mijn L staat op (4,1) en C op (7,1)! En de andere speler zegt dat mijn L op (2,1) staat... maar dat klopt niet. (2,1) is leeg (◌).} \\
\\ 
\texttt{Wacht, de coördinaten: x is kolom 1{-}7, y is rij 1{-}7. De bovenrand toont x van 1 tot 7. Dus mijn L staat op x=4, y=1? Of is de notatie (x,y) = (kolom, rij)?} \\
\\ 
\texttt{De andere speler zei: "Mijn C staat op (3,1), L op (3,2), P op (3,5)". En ik antwoordde dat mijn C op (7,1), L op (2,1), P op (3,5) staat. Maar kijkend naar het bord: L staat op rij 1, en afgaande op de tekst "◌◌│L◌◌C" {-} de L is het vierde teken in de reeks na de rand ║. Dus kolom 4, rij 1 = (4,1). En C is het zevende teken = (7,1).} \\
\\ 
\texttt{Hmm, maar ik heb zelf gezegd dat L op (2,1) staat. Dat was mijn fout. De andere speler baseert zich daarop.} \\
\\ 
\texttt{Ik moet dit corrigeren. Mijn L staat op (4,1), niet (2,1). Het voorstel om L op (3,1) te zetten kan niet, want (3,1) is een muur (│).} \\
\\ 
\texttt{Ik moet de andere speler informeren over mijn fout en de juiste posities geven.} \\
            }
        }
    }
     \\ \\

    \theutterance \stepcounter{utterance}  
    & & & \multicolumn{2}{p{0.3\linewidth}}{
        \cellcolor[rgb]{0.9,0.9,0.9}{
            \makecell[{{p{\linewidth}}}]{
                \texttt{\tiny{[GM$|$GM]}}
                \texttt{Ik analyseer het voorstel van de andere speler:} \\
\\ 
\texttt{{-} Doel: C op (2,1), L op (3,1), P op (3,5)} \\
\texttt{{-} Mijn huidige posities: C op (7,1), L op (2,1), P op (3,5)} \\
\texttt{{-} Mijn L staat al op (2,1), maar volgens het voorstel moet L naar (3,1). Mijn C moet van (7,1) naar (2,1) of (3,1)?} \\
\\ 
\texttt{Wacht, even goed kijken. Het voorstel is C op (2,1) en L op (3,1). Mijn L staat op (2,1) en mijn C op (7,1). Dus ik moet L verplaatsen van (2,1) naar (3,1), en C van (7,1) naar (2,1). Dat klopt.} \\
\\ 
\texttt{Maar ik moet eerst controleren of (2,1) en (3,1) leeg zijn. Mijn L staat op (2,1), dus als ik L verplaats naar (3,1), dan komt (2,1) vrij voor C. (3,1) moet leeg zijn. Laat ik het bord checken: op rij 1 zijn de cellen (1,1) en (2,1) vóór de │. Mijn L staat op (2,1), C op (7,1). Positie (3,1) {-} even kijken naar het bord. De bovenste rij is: ╔══╤════╗ met rij 1: ║◌◌│L◌◌C║. Dus (1,1)=◌, (2,1)=◌ maar daar staat L, (3,1)=│ (muur), (4,1)=L? Nee, (2,1)=L, daarna │, dan (4,1)=◌, (5,1)=◌, (6,1)=C? Even opnieuw.} \\
\\ 
\texttt{Rij 1: ║◌◌│L◌◌C║} \\
\texttt{{-} (1,1) = ◌} \\
\texttt{{-} (2,1) = ◌ maar hier staat L? Nee, "◌◌│L◌◌C" {-} de eerste twee zijn ◌, dan │, dan L, ◌, ◌, C. Dus:} \\
\texttt{{-} (1,1) = ◌} \\
\texttt{{-} (2,1) = ◌} \\
\texttt{{-} (3,1) = │ (muur)} \\
\texttt{{-} (4,1) = L} \\
\texttt{{-} (5,1) = ◌} \\
\texttt{{-} (6,1) = ◌} \\
\texttt{{-} (7,1) = C} \\
\\ 
\texttt{Oh! Mijn L staat op (4,1) en C op (7,1)! En de andere speler zegt dat mijn L op (2,1) staat... maar dat klopt niet. (2,1) is leeg (◌).} \\
\\ 
\texttt{Wacht, de coördinaten: x is kolom 1{-}7, y is rij 1{-}7. De bovenrand toont x van 1 tot 7. Dus mijn L staat op x=4, y=1? Of is de notatie (x,y) = (kolom, rij)?} \\
\\ 
\texttt{De andere speler zei: "Mijn C staat op (3,1), L op (3,2), P op (3,5)". En ik antwoordde dat mijn C op (7,1), L op (2,1), P op (3,5) staat. Maar kijkend naar het bord: L staat op rij 1, en afgaande op de tekst "◌◌│L◌◌C" {-} de L is het vierde teken in de reeks na de rand ║. Dus kolom 4, rij 1 = (4,1). En C is het zevende teken = (7,1).} \\
\\ 
\texttt{Hmm, maar ik heb zelf gezegd dat L op (2,1) staat. Dat was mijn fout. De andere speler baseert zich daarop.} \\
\\ 
\texttt{Ik moet dit corrigeren. Mijn L staat op (4,1), niet (2,1). Het voorstel om L op (3,1) te zetten kan niet, want (3,1) is een muur (│).} \\
\\ 
\texttt{Ik moet de andere speler informeren over mijn fout en de juiste posities geven.} \\
            }
        }
    }
    & & \\ \\

    \theutterance \stepcounter{utterance}  
    & & & \multicolumn{2}{p{0.3\linewidth}}{
        \cellcolor[rgb]{0.9,0.9,0.9}{
            \makecell[{{p{\linewidth}}}]{
                \texttt{\tiny{[GM$|$GM]}}
                \texttt{Invalid response format} \\
            }
        }
    }
    & & \\ \\

    \theutterance \stepcounter{utterance}  
    & & \multicolumn{4}{p{0.6\linewidth}}{
        \cellcolor[rgb]{0.9,0.9,0.9}{
            \makecell[{{p{\linewidth}}}]{
                \texttt{\tiny{[GM$\rangle$P2]}}
                \texttt{Straf: Je bericht heeft niet het verwachte formaat!} \\
\\ 
\texttt{Zorg ervoor dat je antwoord alleen `ZEG: <BERICHT>` of `BEWEEG: <OBJECT>, (<X>, <Y>)` bevat, en niets anders!} \\
\\ 
\texttt{Jullie hebben gezamenlijk 3 van 8 straffen verzameld.} \\
\texttt{Probeer het opnieuw!} \\
            }
        }
    }
    & \\ \\

    \theutterance \stepcounter{utterance}  
    & & & \multicolumn{4}{p{0.6\linewidth}}{
        \cellcolor[rgb]{0.9,0.9,0.9}{
            \makecell[{{p{\linewidth}}}]{
                \texttt{\tiny{[GM$\langle$P2]}}
                \texttt{ZEG: Sorry, ik maakte een fout! Mijn L staat niet op (2,1) maar op (4,1). Mijn C op (7,1), P op (3,5). (2,1) is leeg. En (3,1) is een muur, dus daar kunnen we geen object plaatsen. Wil je opnieuw voorstellen waar we C en L plaatsen?} \\
            }
        }
    }
     \\ \\

    \theutterance \stepcounter{utterance}  
    & & & \multicolumn{2}{p{0.3\linewidth}}{
        \cellcolor[rgb]{0.9,0.9,0.9}{
            \makecell[{{p{\linewidth}}}]{
                \texttt{\tiny{[GM$|$GM]}}
                \texttt{ZEG: Sorry, ik maakte een fout! Mijn L staat niet op (2,1) maar op (4,1). Mijn C op (7,1), P op (3,5). (2,1) is leeg. En (3,1) is een muur, dus daar kunnen we geen object plaatsen. Wil je opnieuw voorstellen waar we C en L plaatsen?} \\
            }
        }
    }
    & & \\ \\

    \theutterance \stepcounter{utterance}  
    & & \multicolumn{4}{p{0.6\linewidth}}{
        \cellcolor[rgb]{0.9,0.9,0.9}{
            \makecell[{{p{\linewidth}}}]{
                \texttt{\tiny{[P1$\langle$GM]}}
                \texttt{Je bericht is doorgestuurd naar de andere speler.} \\
\\ 
\\ 
\texttt{Je speelt momenteel ronde 2 van maximaal 12.} \\
\texttt{Jullie hebben gezamenlijk 3 van 8 straffen verzameld.} \\
\texttt{De andere speler stuurde dit bericht:} \\
\\ 
\\ 
\texttt{"Sorry, ik maakte een fout! Mijn L staat niet op (2,1) maar op (4,1). Mijn C op (7,1), P op (3,5). (2,1) is leeg. En (3,1) is een muur, dus daar kunnen we geen object plaatsen. Wil je opnieuw voorstellen waar we C en L plaatsen?"} \\
\\ 
\\ 
\texttt{Wat is je volgende opdracht?} \\
            }
        }
    }
    & \\ \\

    \theutterance \stepcounter{utterance}  
    & \multicolumn{4}{p{0.6\linewidth}}{
        \cellcolor[rgb]{0.9,0.9,0.9}{
            \makecell[{{p{\linewidth}}}]{
                \texttt{\tiny{[P1$\rangle$GM]}}
                \texttt{ZEG: Geen probleem! (3,1) is op mijn bord leeg, maar als dat bij jou een muur is moeten we iets anders kiezen. Laten we ons richten op lege cellen die voor ons beiden beschikbaar zijn. Mijn C staat nu op (3,1), die kan ik verplaatsen. Mijn L staat op (3,2). Stel voor: C naar (6,1) en L naar (4,2)? Of heb jij daar muren? Noem anders twee lege cellen die bij jou beschikbaar zijn.} \\
            }
        }
    }
    & & \\ \\

    \theutterance \stepcounter{utterance}  
    & & & \multicolumn{2}{p{0.3\linewidth}}{
        \cellcolor[rgb]{0.9,0.9,0.9}{
            \makecell[{{p{\linewidth}}}]{
                \texttt{\tiny{[GM$|$GM]}}
                \texttt{ZEG: Geen probleem! (3,1) is op mijn bord leeg, maar als dat bij jou een muur is moeten we iets anders kiezen. Laten we ons richten op lege cellen die voor ons beiden beschikbaar zijn. Mijn C staat nu op (3,1), die kan ik verplaatsen. Mijn L staat op (3,2). Stel voor: C naar (6,1) en L naar (4,2)? Of heb jij daar muren? Noem anders twee lege cellen die bij jou beschikbaar zijn.} \\
            }
        }
    }
    & & \\ \\

    \theutterance \stepcounter{utterance}  
    & & \multicolumn{4}{p{0.6\linewidth}}{
        \cellcolor[rgb]{0.9,0.9,0.9}{
            \makecell[{{p{\linewidth}}}]{
                \texttt{\tiny{[GM$\rangle$P2]}}
                \texttt{Je bericht is doorgestuurd naar de andere speler.} \\
\\ 
\\ 
\texttt{Je speelt momenteel ronde 3 van maximaal 12.} \\
\texttt{Jullie hebben gezamenlijk 3 van 8 straffen verzameld.} \\
\texttt{De andere speler stuurde dit bericht:} \\
\\ 
\\ 
\texttt{"Geen probleem! (3,1) is op mijn bord leeg, maar als dat bij jou een muur is moeten we iets anders kiezen. Laten we ons richten op lege cellen die voor ons beiden beschikbaar zijn. Mijn C staat nu op (3,1), die kan ik verplaatsen. Mijn L staat op (3,2). Stel voor: C naar (6,1) en L naar (4,2)? Of heb jij daar muren? Noem anders twee lege cellen die bij jou beschikbaar zijn."} \\
\\ 
\\ 
\texttt{Wat is je volgende opdracht?} \\
            }
        }
    }
    & \\ \\

    \theutterance \stepcounter{utterance}  
    & & & \multicolumn{4}{p{0.6\linewidth}}{
        \cellcolor[rgb]{0.9,0.9,0.9}{
            \makecell[{{p{\linewidth}}}]{
                \texttt{\tiny{[GM$\langle$P2]}}
                \texttt{Ik moet controleren welke cellen op mijn bord leeg zijn en geen muren bevatten. Laat ik mijn bord nog eens goed bekijken:} \\
\\ 
\texttt{Rij 1: ║◌◌│L◌◌C║ → (1,1)=◌, (2,1)=◌, (3,1)=│, (4,1)=L, (5,1)=◌, (6,1)=◌, (7,1)=C} \\
\texttt{Rij 2: ╟──┤◌◌┌─╢ → (1,2)=─, (2,2)=─, (3,2)=┤, (4,2)=◌, (5,2)=◌, (6,2)=┌, (7,2)=─} \\
\texttt{Rij 3: ║◌◌├──┤◌║ → (1,3)=◌, (2,3)=◌, (3,3)=├, (4,3)=─, (5,3)=─, (6,3)=┤, (7,3)=◌} \\
\texttt{Rij 4: ╟──┘◌◌├─╢ → (1,4)=─, (2,4)=─, (3,4)=┘, (4,4)=◌, (5,4)=◌, (6,4)=├, (7,4)=─} \\
\texttt{Rij 5: ║◌◌P◌◌│◌║ → (1,5)=◌, (2,5)=◌, (3,5)=P, (4,5)=◌, (5,5)=◌, (6,5)=│, (7,5)=◌} \\
\texttt{Rij 6: ╟─┐◌┌─┼─╢ → (1,6)=─, (2,6)=┐, (3,6)=◌, (4,6)=┌, (5,6)=─, (6,6)=┼, (7,6)=─} \\
\texttt{Rij 7: ║◌│◌│◌│◌║ → (1,7)=◌, (2,7)=│, (3,7)=◌, (4,7)=│, (5,7)=◌, (6,7)=│, (7,7)=◌} \\
\\ 
\texttt{Lege cellen (◌) op mijn bord:} \\
\texttt{(1,1), (2,1), (5,1), (6,1)} \\
\texttt{(4,2), (5,2)} \\
\texttt{(1,3), (2,3), (7,3)} \\
\texttt{(4,4), (5,4)} \\
\texttt{(1,5), (2,5), (4,5), (5,5), (7,5)} \\
\texttt{(3,6)} \\
\texttt{(1,7), (3,7), (5,7), (7,7)} \\
\\ 
\texttt{Het voorstel van de ander: C naar (6,1) en L naar (4,2). (6,1) is leeg op mijn bord, en (4,2) is ook leeg. Dat lijkt te werken! Mijn C staat op (7,1) en moet naar (6,1). Mijn L staat op (4,1) en moet naar (4,2). P staat al op (3,5) bij ons beiden. Dit klinkt goed!ZEG: (6,1) en (4,2) zijn bij mij leeg en beschikbaar! Dat voorstel werkt perfect. Ik verplaats mijn C van (7,1) naar (6,1) en mijn L van (4,1) naar (4,2). Zullen we dat als doel kiezen? P blijft op (3,5).} \\
            }
        }
    }
     \\ \\

    \theutterance \stepcounter{utterance}  
    & & & \multicolumn{2}{p{0.3\linewidth}}{
        \cellcolor[rgb]{0.9,0.9,0.9}{
            \makecell[{{p{\linewidth}}}]{
                \texttt{\tiny{[GM$|$GM]}}
                \texttt{Ik moet controleren welke cellen op mijn bord leeg zijn en geen muren bevatten. Laat ik mijn bord nog eens goed bekijken:} \\
\\ 
\texttt{Rij 1: ║◌◌│L◌◌C║ → (1,1)=◌, (2,1)=◌, (3,1)=│, (4,1)=L, (5,1)=◌, (6,1)=◌, (7,1)=C} \\
\texttt{Rij 2: ╟──┤◌◌┌─╢ → (1,2)=─, (2,2)=─, (3,2)=┤, (4,2)=◌, (5,2)=◌, (6,2)=┌, (7,2)=─} \\
\texttt{Rij 3: ║◌◌├──┤◌║ → (1,3)=◌, (2,3)=◌, (3,3)=├, (4,3)=─, (5,3)=─, (6,3)=┤, (7,3)=◌} \\
\texttt{Rij 4: ╟──┘◌◌├─╢ → (1,4)=─, (2,4)=─, (3,4)=┘, (4,4)=◌, (5,4)=◌, (6,4)=├, (7,4)=─} \\
\texttt{Rij 5: ║◌◌P◌◌│◌║ → (1,5)=◌, (2,5)=◌, (3,5)=P, (4,5)=◌, (5,5)=◌, (6,5)=│, (7,5)=◌} \\
\texttt{Rij 6: ╟─┐◌┌─┼─╢ → (1,6)=─, (2,6)=┐, (3,6)=◌, (4,6)=┌, (5,6)=─, (6,6)=┼, (7,6)=─} \\
\texttt{Rij 7: ║◌│◌│◌│◌║ → (1,7)=◌, (2,7)=│, (3,7)=◌, (4,7)=│, (5,7)=◌, (6,7)=│, (7,7)=◌} \\
\\ 
\texttt{Lege cellen (◌) op mijn bord:} \\
\texttt{(1,1), (2,1), (5,1), (6,1)} \\
\texttt{(4,2), (5,2)} \\
\texttt{(1,3), (2,3), (7,3)} \\
\texttt{(4,4), (5,4)} \\
\texttt{(1,5), (2,5), (4,5), (5,5), (7,5)} \\
\texttt{(3,6)} \\
\texttt{(1,7), (3,7), (5,7), (7,7)} \\
\\ 
\texttt{Het voorstel van de ander: C naar (6,1) en L naar (4,2). (6,1) is leeg op mijn bord, en (4,2) is ook leeg. Dat lijkt te werken! Mijn C staat op (7,1) en moet naar (6,1). Mijn L staat op (4,1) en moet naar (4,2). P staat al op (3,5) bij ons beiden. Dit klinkt goed!ZEG: (6,1) en (4,2) zijn bij mij leeg en beschikbaar! Dat voorstel werkt perfect. Ik verplaats mijn C van (7,1) naar (6,1) en mijn L van (4,1) naar (4,2). Zullen we dat als doel kiezen? P blijft op (3,5).} \\
            }
        }
    }
    & & \\ \\

    \theutterance \stepcounter{utterance}  
    & & & \multicolumn{2}{p{0.3\linewidth}}{
        \cellcolor[rgb]{0.9,0.9,0.9}{
            \makecell[{{p{\linewidth}}}]{
                \texttt{\tiny{[GM$|$GM]}}
                \texttt{Invalid format: head is not empty: 'Ik moet controleren welke cellen op mijn bord leeg zijn en geen muren bevatten. Laat ik mijn bord nog eens goed bekijken:} \\
\\ 
\texttt{Rij 1: ║◌◌│L◌◌C║ → (1,1)=◌, (2,1)=◌, (3,1)=│, (4,1)=L, (5,1)=◌, (6,1)=◌, (7,1)=C} \\
\texttt{Rij 2: ╟──┤◌◌┌─╢ → (1,2)=─, (2,2)=─, (3,2)=┤, (4,2)=◌, (5,2)=◌, (6,2)=┌, (7,2)=─} \\
\texttt{Rij 3: ║◌◌├──┤◌║ → (1,3)=◌, (2,3)=◌, (3,3)=├, (4,3)=─, (5,3)=─, (6,3)=┤, (7,3)=◌} \\
\texttt{Rij 4: ╟──┘◌◌├─╢ → (1,4)=─, (2,4)=─, (3,4)=┘, (4,4)=◌, (5,4)=◌, (6,4)=├, (7,4)=─} \\
\texttt{Rij 5: ║◌◌P◌◌│◌║ → (1,5)=◌, (2,5)=◌, (3,5)=P, (4,5)=◌, (5,5)=◌, (6,5)=│, (7,5)=◌} \\
\texttt{Rij 6: ╟─┐◌┌─┼─╢ → (1,6)=─, (2,6)=┐, (3,6)=◌, (4,6)=┌, (5,6)=─, (6,6)=┼, (7,6)=─} \\
\texttt{Rij 7: ║◌│◌│◌│◌║ → (1,7)=◌, (2,7)=│, (3,7)=◌, (4,7)=│, (5,7)=◌, (6,7)=│, (7,7)=◌} \\
\\ 
\texttt{Lege cellen (◌) op mijn bord:} \\
\texttt{(1,1), (2,1), (5,1), (6,1)} \\
\texttt{(4,2), (5,2)} \\
\texttt{(1,3), (2,3), (7,3)} \\
\texttt{(4,4), (5,4)} \\
\texttt{(1,5), (2,5), (4,5), (5,5), (7,5)} \\
\texttt{(3,6)} \\
\texttt{(1,7), (3,7), (5,7), (7,7)} \\
\\ 
\texttt{Het voorstel van de ander: C naar (6,1) en L naar (4,2). (6,1) is leeg op mijn bord, en (4,2) is ook leeg. Dat lijkt te werken! Mijn C staat op (7,1) en moet naar (6,1). Mijn L staat op (4,1) en moet naar (4,2). P staat al op (3,5) bij ons beiden. Dit klinkt goed!'} \\
            }
        }
    }
    & & \\ \\

    \theutterance \stepcounter{utterance}  
    & & \multicolumn{4}{p{0.6\linewidth}}{
        \cellcolor[rgb]{0.9,0.9,0.9}{
            \makecell[{{p{\linewidth}}}]{
                \texttt{\tiny{[GM$\rangle$P2]}}
                \texttt{Straf: Je bericht mag niets vóór het commando bevatten!} \\
\\ 
\texttt{Zorg ervoor dat je antwoord alleen `ZEG: <BERICHT>` of `BEWEEG: <OBJECT>, (<X>, <Y>)` bevat, en niets anders!} \\
\\ 
\texttt{Jullie hebben gezamenlijk 4 van 8 straffen verzameld.} \\
\texttt{Probeer het opnieuw!} \\
            }
        }
    }
    & \\ \\

    \theutterance \stepcounter{utterance}  
    & & & \multicolumn{4}{p{0.6\linewidth}}{
        \cellcolor[rgb]{0.9,0.9,0.9}{
            \makecell[{{p{\linewidth}}}]{
                \texttt{\tiny{[GM$\langle$P2]}}
                \texttt{ZEG: (6,1) en (4,2) zijn bij mij leeg en beschikbaar! Dat voorstel werkt perfect. Ik verplaats mijn C van (7,1) naar (6,1) en mijn L van (4,1) naar (4,2). Zullen we dat als doel kiezen? P blijft op (3,5).} \\
            }
        }
    }
     \\ \\

    \theutterance \stepcounter{utterance}  
    & & & \multicolumn{2}{p{0.3\linewidth}}{
        \cellcolor[rgb]{0.9,0.9,0.9}{
            \makecell[{{p{\linewidth}}}]{
                \texttt{\tiny{[GM$|$GM]}}
                \texttt{ZEG: (6,1) en (4,2) zijn bij mij leeg en beschikbaar! Dat voorstel werkt perfect. Ik verplaats mijn C van (7,1) naar (6,1) en mijn L van (4,1) naar (4,2). Zullen we dat als doel kiezen? P blijft op (3,5).} \\
            }
        }
    }
    & & \\ \\

    \theutterance \stepcounter{utterance}  
    & & \multicolumn{4}{p{0.6\linewidth}}{
        \cellcolor[rgb]{0.9,0.9,0.9}{
            \makecell[{{p{\linewidth}}}]{
                \texttt{\tiny{[P1$\langle$GM]}}
                \texttt{Je bericht is doorgestuurd naar de andere speler.} \\
\\ 
\\ 
\texttt{Je speelt momenteel ronde 3 van maximaal 12.} \\
\texttt{Jullie hebben gezamenlijk 4 van 8 straffen verzameld.} \\
\texttt{De andere speler stuurde dit bericht:} \\
\\ 
\\ 
\texttt{"(6,1) en (4,2) zijn bij mij leeg en beschikbaar! Dat voorstel werkt perfect. Ik verplaats mijn C van (7,1) naar (6,1) en mijn L van (4,1) naar (4,2). Zullen we dat als doel kiezen? P blijft op (3,5)."} \\
\\ 
\\ 
\texttt{Wat is je volgende opdracht?} \\
            }
        }
    }
    & \\ \\

    \theutterance \stepcounter{utterance}  
    & \multicolumn{4}{p{0.6\linewidth}}{
        \cellcolor[rgb]{0.9,0.9,0.9}{
            \makecell[{{p{\linewidth}}}]{
                \texttt{\tiny{[P1$\rangle$GM]}}
                \texttt{ZEG: Perfect! Doel: C op (6,1), L op (4,2), P op (3,5). Ik ga nu mijn C van (3,1) naar (6,1) verplaatsen en mijn L van (3,2) naar (4,2). Succes!} \\
            }
        }
    }
    & & \\ \\

    \theutterance \stepcounter{utterance}  
    & & & \multicolumn{2}{p{0.3\linewidth}}{
        \cellcolor[rgb]{0.9,0.9,0.9}{
            \makecell[{{p{\linewidth}}}]{
                \texttt{\tiny{[GM$|$GM]}}
                \texttt{ZEG: Perfect! Doel: C op (6,1), L op (4,2), P op (3,5). Ik ga nu mijn C van (3,1) naar (6,1) verplaatsen en mijn L van (3,2) naar (4,2). Succes!} \\
            }
        }
    }
    & & \\ \\

    \theutterance \stepcounter{utterance}  
    & & \multicolumn{4}{p{0.6\linewidth}}{
        \cellcolor[rgb]{0.9,0.9,0.9}{
            \makecell[{{p{\linewidth}}}]{
                \texttt{\tiny{[GM$\rangle$P2]}}
                \texttt{Je bericht is doorgestuurd naar de andere speler.} \\
\\ 
\\ 
\texttt{Je speelt momenteel ronde 4 van maximaal 12.} \\
\texttt{Jullie hebben gezamenlijk 4 van 8 straffen verzameld.} \\
\texttt{De andere speler stuurde dit bericht:} \\
\\ 
\\ 
\texttt{"Perfect! Doel: C op (6,1), L op (4,2), P op (3,5). Ik ga nu mijn C van (3,1) naar (6,1) verplaatsen en mijn L van (3,2) naar (4,2). Succes!"} \\
\\ 
\\ 
\texttt{Wat is je volgende opdracht?} \\
            }
        }
    }
    & \\ \\

    \theutterance \stepcounter{utterance}  
    & & & \multicolumn{4}{p{0.6\linewidth}}{
        \cellcolor[rgb]{0.9,0.9,0.9}{
            \makecell[{{p{\linewidth}}}]{
                \texttt{\tiny{[GM$\langle$P2]}}
                \texttt{De andere speler is akkoord. Nu ga ik mijn zetten doen. Eerst verplaats ik mijn C van (7,1) naar (6,1). (6,1) is leeg op mijn bord, dus dat zou moeten lukken.} \\
\\ 
\texttt{BEWEEG: C, (6, 1)} \\
            }
        }
    }
     \\ \\

    \theutterance \stepcounter{utterance}  
    & & & \multicolumn{2}{p{0.3\linewidth}}{
        \cellcolor[rgb]{0.9,0.9,0.9}{
            \makecell[{{p{\linewidth}}}]{
                \texttt{\tiny{[GM$|$GM]}}
                \texttt{De andere speler is akkoord. Nu ga ik mijn zetten doen. Eerst verplaats ik mijn C van (7,1) naar (6,1). (6,1) is leeg op mijn bord, dus dat zou moeten lukken.} \\
\\ 
\texttt{BEWEEG: C, (6, 1)} \\
            }
        }
    }
    & & \\ \\

    \theutterance \stepcounter{utterance}  
    & & & \multicolumn{2}{p{0.3\linewidth}}{
        \cellcolor[rgb]{0.9,0.9,0.9}{
            \makecell[{{p{\linewidth}}}]{
                \texttt{\tiny{[GM$|$GM]}}
                \texttt{Invalid format: head is not empty: 'De andere speler is akkoord. Nu ga ik mijn zetten doen. Eerst verplaats ik mijn C van (7,1) naar (6,1). (6,1) is leeg op mijn bord, dus dat zou moeten lukken.} \\
\\ 
\texttt{'} \\
            }
        }
    }
    & & \\ \\

    \theutterance \stepcounter{utterance}  
    & & \multicolumn{4}{p{0.6\linewidth}}{
        \cellcolor[rgb]{0.9,0.9,0.9}{
            \makecell[{{p{\linewidth}}}]{
                \texttt{\tiny{[GM$\rangle$P2]}}
                \texttt{Straf: Je bericht mag niets vóór het commando bevatten!} \\
\\ 
\texttt{Zorg ervoor dat je antwoord alleen `ZEG: <BERICHT>` of `BEWEEG: <OBJECT>, (<X>, <Y>)` bevat, en niets anders!} \\
\\ 
\texttt{Jullie hebben gezamenlijk 5 van 8 straffen verzameld.} \\
\texttt{Probeer het opnieuw!} \\
            }
        }
    }
    & \\ \\

    \theutterance \stepcounter{utterance}  
    & & & \multicolumn{4}{p{0.6\linewidth}}{
        \cellcolor[rgb]{0.9,0.9,0.9}{
            \makecell[{{p{\linewidth}}}]{
                \texttt{\tiny{[GM$\langle$P2]}}
                \texttt{BEWEEG: C, (6, 1)} \\
            }
        }
    }
     \\ \\

    \theutterance \stepcounter{utterance}  
    & & & \multicolumn{2}{p{0.3\linewidth}}{
        \cellcolor[rgb]{0.9,0.9,0.9}{
            \makecell[{{p{\linewidth}}}]{
                \texttt{\tiny{[GM$|$GM]}}
                \texttt{BEWEEG: C, (6, 1)} \\
            }
        }
    }
    & & \\ \\

    \theutterance \stepcounter{utterance}  
    & & & \multicolumn{2}{p{0.3\linewidth}}{
        \cellcolor[rgb]{0.9,0.9,0.9}{
            \makecell[{{p{\linewidth}}}]{
                \texttt{\tiny{[GM$|$GM]}}
                \texttt{C moved to (6, 1)} \\
            }
        }
    }
    & & \\ \\

    \theutterance \stepcounter{utterance}  
    & & \multicolumn{4}{p{0.6\linewidth}}{
        \cellcolor[rgb]{0.9,0.9,0.9}{
            \makecell[{{p{\linewidth}}}]{
                \texttt{\tiny{[P1$\langle$GM]}}
                \texttt{Je bericht is doorgestuurd naar de andere speler.} \\
\\ 
\\ 
\texttt{Je speelt momenteel ronde 4 van maximaal 12.} \\
\texttt{Jullie hebben gezamenlijk 5 van 8 straffen verzameld.} \\
\texttt{De andere speler heeft een voorwerp op zijn rooster verplaatst.} \\
\\ 
\\ 
\texttt{Wat is je volgende opdracht?} \\
            }
        }
    }
    & \\ \\

    \theutterance \stepcounter{utterance}  
    & \multicolumn{4}{p{0.6\linewidth}}{
        \cellcolor[rgb]{0.9,0.9,0.9}{
            \makecell[{{p{\linewidth}}}]{
                \texttt{\tiny{[P1$\rangle$GM]}}
                \texttt{BEWEEG: C, (6, 1)} \\
            }
        }
    }
    & & \\ \\

    \theutterance \stepcounter{utterance}  
    & & & \multicolumn{2}{p{0.3\linewidth}}{
        \cellcolor[rgb]{0.9,0.9,0.9}{
            \makecell[{{p{\linewidth}}}]{
                \texttt{\tiny{[GM$|$GM]}}
                \texttt{BEWEEG: C, (6, 1)} \\
            }
        }
    }
    & & \\ \\

    \theutterance \stepcounter{utterance}  
    & & & \multicolumn{2}{p{0.3\linewidth}}{
        \cellcolor[rgb]{0.9,0.9,0.9}{
            \makecell[{{p{\linewidth}}}]{
                \texttt{\tiny{[GM$|$GM]}}
                \texttt{C moved to (6, 1)} \\
            }
        }
    }
    & & \\ \\

    \theutterance \stepcounter{utterance}  
    & & \multicolumn{4}{p{0.6\linewidth}}{
        \cellcolor[rgb]{0.9,0.9,0.9}{
            \makecell[{{p{\linewidth}}}]{
                \texttt{\tiny{[GM$\rangle$P2]}}
                \texttt{Verplaatst 'C' naar (6,1) met succes. Uw bijgewerkte raster ziet er als volgt uit:} \\
\texttt{```} \\
\texttt{ 1234567} \\
\texttt{╔══╤════╗} \\
\texttt{║◌◌│L◌C◌║ 1} \\
\texttt{╟──┤◌◌┌─╢ 2} \\
\texttt{║◌◌├──┤◌║ 3} \\
\texttt{╟──┘◌◌├─╢ 4} \\
\texttt{║◌◌P◌◌│◌║ 5} \\
\texttt{╟─┐◌┌─┼─╢ 6} \\
\texttt{║◌│◌│◌│◌║ 7} \\
\texttt{╚═╧═╧═╧═╝} \\
\texttt{```} \\
\texttt{Je speelt momenteel ronde 5 van maximaal 12.} \\
\texttt{Jullie hebben gezamenlijk 5 van 8 straffen verzameld.} \\
\texttt{De andere speler heeft een voorwerp op zijn rooster verplaatst.} \\
\\ 
\\ 
\texttt{Wat is je volgende opdracht?} \\
            }
        }
    }
    & \\ \\

    \theutterance \stepcounter{utterance}  
    & & & \multicolumn{4}{p{0.6\linewidth}}{
        \cellcolor[rgb]{0.9,0.9,0.9}{
            \makecell[{{p{\linewidth}}}]{
                \texttt{\tiny{[GM$\langle$P2]}}
                \texttt{Nu verplaats ik mijn L van (4,1) naar (4,2). Op mijn bord is (4,2) leeg (◌), dus dat zou moeten lukken.} \\
\\ 
\texttt{BEWEEG: L, (4, 2)} \\
            }
        }
    }
     \\ \\

    \theutterance \stepcounter{utterance}  
    & & & \multicolumn{2}{p{0.3\linewidth}}{
        \cellcolor[rgb]{0.9,0.9,0.9}{
            \makecell[{{p{\linewidth}}}]{
                \texttt{\tiny{[GM$|$GM]}}
                \texttt{Nu verplaats ik mijn L van (4,1) naar (4,2). Op mijn bord is (4,2) leeg (◌), dus dat zou moeten lukken.} \\
\\ 
\texttt{BEWEEG: L, (4, 2)} \\
            }
        }
    }
    & & \\ \\

    \theutterance \stepcounter{utterance}  
    & & & \multicolumn{2}{p{0.3\linewidth}}{
        \cellcolor[rgb]{0.9,0.9,0.9}{
            \makecell[{{p{\linewidth}}}]{
                \texttt{\tiny{[GM$|$GM]}}
                \texttt{Invalid format: head is not empty: 'Nu verplaats ik mijn L van (4,1) naar (4,2). Op mijn bord is (4,2) leeg (◌), dus dat zou moeten lukken.} \\
\\ 
\texttt{'} \\
            }
        }
    }
    & & \\ \\

    \theutterance \stepcounter{utterance}  
    & & \multicolumn{4}{p{0.6\linewidth}}{
        \cellcolor[rgb]{0.9,0.9,0.9}{
            \makecell[{{p{\linewidth}}}]{
                \texttt{\tiny{[GM$\rangle$P2]}}
                \texttt{Straf: Je bericht mag niets vóór het commando bevatten!} \\
\\ 
\texttt{Zorg ervoor dat je antwoord alleen `ZEG: <BERICHT>` of `BEWEEG: <OBJECT>, (<X>, <Y>)` bevat, en niets anders!} \\
\\ 
\texttt{Jullie hebben gezamenlijk 6 van 8 straffen verzameld.} \\
\texttt{Probeer het opnieuw!} \\
            }
        }
    }
    & \\ \\

    \theutterance \stepcounter{utterance}  
    & & & \multicolumn{4}{p{0.6\linewidth}}{
        \cellcolor[rgb]{0.9,0.9,0.9}{
            \makecell[{{p{\linewidth}}}]{
                \texttt{\tiny{[GM$\langle$P2]}}
                \texttt{BEWEEG: L, (4, 2)} \\
            }
        }
    }
     \\ \\

    \theutterance \stepcounter{utterance}  
    & & & \multicolumn{2}{p{0.3\linewidth}}{
        \cellcolor[rgb]{0.9,0.9,0.9}{
            \makecell[{{p{\linewidth}}}]{
                \texttt{\tiny{[GM$|$GM]}}
                \texttt{BEWEEG: L, (4, 2)} \\
            }
        }
    }
    & & \\ \\

    \theutterance \stepcounter{utterance}  
    & & & \multicolumn{2}{p{0.3\linewidth}}{
        \cellcolor[rgb]{0.9,0.9,0.9}{
            \makecell[{{p{\linewidth}}}]{
                \texttt{\tiny{[GM$|$GM]}}
                \texttt{L moved to (4, 2)} \\
            }
        }
    }
    & & \\ \\

    \theutterance \stepcounter{utterance}  
    & & \multicolumn{4}{p{0.6\linewidth}}{
        \cellcolor[rgb]{0.9,0.9,0.9}{
            \makecell[{{p{\linewidth}}}]{
                \texttt{\tiny{[P1$\langle$GM]}}
                \texttt{Verplaatst 'C' naar (6,1) met succes. Uw bijgewerkte raster ziet er als volgt uit:} \\
\texttt{```} \\
\texttt{ 1234567} \\
\texttt{╔══╤════╗} \\
\texttt{║◌◌│◌◌C◌║ 1} \\
\texttt{╟──┤◌L┌─╢ 2} \\
\texttt{║◌◌├──┤◌║ 3} \\
\texttt{╟──┘◌◌├─╢ 4} \\
\texttt{║◌◌P◌◌│◌║ 5} \\
\texttt{╟─┐◌┌─┼─╢ 6} \\
\texttt{║◌│◌│◌│◌║ 7} \\
\texttt{╚═╧═╧═╧═╝} \\
\texttt{```} \\
\texttt{Je speelt momenteel ronde 5 van maximaal 12.} \\
\texttt{Jullie hebben gezamenlijk 6 van 8 straffen verzameld.} \\
\texttt{De andere speler heeft een voorwerp op zijn rooster verplaatst.} \\
\\ 
\\ 
\texttt{Wat is je volgende opdracht?} \\
            }
        }
    }
    & \\ \\

    \theutterance \stepcounter{utterance}  
    & \multicolumn{4}{p{0.6\linewidth}}{
        \cellcolor[rgb]{0.9,0.9,0.9}{
            \makecell[{{p{\linewidth}}}]{
                \texttt{\tiny{[P1$\rangle$GM]}}
                \texttt{BEWEEG: L, (4, 2)} \\
            }
        }
    }
    & & \\ \\

    \theutterance \stepcounter{utterance}  
    & & & \multicolumn{2}{p{0.3\linewidth}}{
        \cellcolor[rgb]{0.9,0.9,0.9}{
            \makecell[{{p{\linewidth}}}]{
                \texttt{\tiny{[GM$|$GM]}}
                \texttt{BEWEEG: L, (4, 2)} \\
            }
        }
    }
    & & \\ \\

    \theutterance \stepcounter{utterance}  
    & & & \multicolumn{2}{p{0.3\linewidth}}{
        \cellcolor[rgb]{0.9,0.9,0.9}{
            \makecell[{{p{\linewidth}}}]{
                \texttt{\tiny{[GM$|$GM]}}
                \texttt{L moved to (4, 2)} \\
            }
        }
    }
    & & \\ \\

    \theutterance \stepcounter{utterance}  
    & & \multicolumn{4}{p{0.6\linewidth}}{
        \cellcolor[rgb]{0.9,0.9,0.9}{
            \makecell[{{p{\linewidth}}}]{
                \texttt{\tiny{[GM$\rangle$P2]}}
                \texttt{Verplaatst 'L' naar (4,2) met succes. Uw bijgewerkte raster ziet er als volgt uit:} \\
\texttt{```} \\
\texttt{ 1234567} \\
\texttt{╔══╤════╗} \\
\texttt{║◌◌│◌◌C◌║ 1} \\
\texttt{╟──┤L◌┌─╢ 2} \\
\texttt{║◌◌├──┤◌║ 3} \\
\texttt{╟──┘◌◌├─╢ 4} \\
\texttt{║◌◌P◌◌│◌║ 5} \\
\texttt{╟─┐◌┌─┼─╢ 6} \\
\texttt{║◌│◌│◌│◌║ 7} \\
\texttt{╚═╧═╧═╧═╝} \\
\texttt{```} \\
\texttt{Je speelt momenteel ronde 6 van maximaal 12.} \\
\texttt{Jullie hebben gezamenlijk 6 van 8 straffen verzameld.} \\
\texttt{De andere speler heeft een voorwerp op zijn rooster verplaatst.} \\
\\ 
\\ 
\texttt{Wat is je volgende opdracht?} \\
            }
        }
    }
    & \\ \\

    \theutterance \stepcounter{utterance}  
    & & & \multicolumn{4}{p{0.6\linewidth}}{
        \cellcolor[rgb]{0.9,0.9,0.9}{
            \makecell[{{p{\linewidth}}}]{
                \texttt{\tiny{[GM$\langle$P2]}}
                \texttt{Mijn objecten staan nu op: C op (6,1), L op (4,2), P op (3,5). Dat komt overeen met het afgesproken doel. Ik wacht even af of de andere speler ook klaar is, en dan kan ik vragen of we klaar zijn.} \\
\\ 
\texttt{ZEG: Ik heb mijn C naar (6,1) en L naar (4,2) verplaatst. Mijn objecten staan nu op de afgesproken posities. Ben jij ook klaar? Klaar?} \\
            }
        }
    }
     \\ \\

    \theutterance \stepcounter{utterance}  
    & & & \multicolumn{2}{p{0.3\linewidth}}{
        \cellcolor[rgb]{0.9,0.9,0.9}{
            \makecell[{{p{\linewidth}}}]{
                \texttt{\tiny{[GM$|$GM]}}
                \texttt{Mijn objecten staan nu op: C op (6,1), L op (4,2), P op (3,5). Dat komt overeen met het afgesproken doel. Ik wacht even af of de andere speler ook klaar is, en dan kan ik vragen of we klaar zijn.} \\
\\ 
\texttt{ZEG: Ik heb mijn C naar (6,1) en L naar (4,2) verplaatst. Mijn objecten staan nu op de afgesproken posities. Ben jij ook klaar? Klaar?} \\
            }
        }
    }
    & & \\ \\

    \theutterance \stepcounter{utterance}  
    & & & \multicolumn{2}{p{0.3\linewidth}}{
        \cellcolor[rgb]{0.9,0.9,0.9}{
            \makecell[{{p{\linewidth}}}]{
                \texttt{\tiny{[GM$|$GM]}}
                \texttt{Invalid format: head is not empty: 'Mijn objecten staan nu op: C op (6,1), L op (4,2), P op (3,5). Dat komt overeen met het afgesproken doel. Ik wacht even af of de andere speler ook klaar is, en dan kan ik vragen of we klaar zijn.} \\
\\ 
\texttt{'} \\
            }
        }
    }
    & & \\ \\

    \theutterance \stepcounter{utterance}  
    & & \multicolumn{4}{p{0.6\linewidth}}{
        \cellcolor[rgb]{0.9,0.9,0.9}{
            \makecell[{{p{\linewidth}}}]{
                \texttt{\tiny{[GM$\rangle$P2]}}
                \texttt{Straf: Je bericht mag niets vóór het commando bevatten!} \\
\\ 
\texttt{Zorg ervoor dat je antwoord alleen `ZEG: <BERICHT>` of `BEWEEG: <OBJECT>, (<X>, <Y>)` bevat, en niets anders!} \\
\\ 
\texttt{Jullie hebben gezamenlijk 7 van 8 straffen verzameld.} \\
\texttt{Probeer het opnieuw!} \\
            }
        }
    }
    & \\ \\

    \theutterance \stepcounter{utterance}  
    & & & \multicolumn{4}{p{0.6\linewidth}}{
        \cellcolor[rgb]{0.9,0.9,0.9}{
            \makecell[{{p{\linewidth}}}]{
                \texttt{\tiny{[GM$\langle$P2]}}
                \texttt{ZEG: Ik heb mijn C naar (6,1) en L naar (4,2) verplaatst. Mijn objecten staan nu op de afgesproken posities. Ben jij ook klaar? Klaar?} \\
            }
        }
    }
     \\ \\

    \theutterance \stepcounter{utterance}  
    & & & \multicolumn{2}{p{0.3\linewidth}}{
        \cellcolor[rgb]{0.9,0.9,0.9}{
            \makecell[{{p{\linewidth}}}]{
                \texttt{\tiny{[GM$|$GM]}}
                \texttt{ZEG: Ik heb mijn C naar (6,1) en L naar (4,2) verplaatst. Mijn objecten staan nu op de afgesproken posities. Ben jij ook klaar? Klaar?} \\
            }
        }
    }
    & & \\ \\

    \theutterance \stepcounter{utterance}  
    & & \multicolumn{4}{p{0.6\linewidth}}{
        \cellcolor[rgb]{0.9,0.9,0.9}{
            \makecell[{{p{\linewidth}}}]{
                \texttt{\tiny{[P1$\langle$GM]}}
                \texttt{Verplaatst 'L' naar (4,2) met succes. Uw bijgewerkte raster ziet er als volgt uit:} \\
\texttt{```} \\
\texttt{ 1234567} \\
\texttt{╔══╤════╗} \\
\texttt{║◌◌│◌◌C◌║ 1} \\
\texttt{╟──┤L◌┌─╢ 2} \\
\texttt{║◌◌├──┤◌║ 3} \\
\texttt{╟──┘◌◌├─╢ 4} \\
\texttt{║◌◌P◌◌│◌║ 5} \\
\texttt{╟─┐◌┌─┼─╢ 6} \\
\texttt{║◌│◌│◌│◌║ 7} \\
\texttt{╚═╧═╧═╧═╝} \\
\texttt{```} \\
\texttt{Je speelt momenteel ronde 6 van maximaal 12.} \\
\texttt{Jullie hebben gezamenlijk 7 van 8 straffen verzameld.} \\
\texttt{De andere speler stuurde dit bericht:} \\
\\ 
\\ 
\texttt{"Ik heb mijn C naar (6,1) en L naar (4,2) verplaatst. Mijn objecten staan nu op de afgesproken posities. Ben jij ook klaar? Klaar?"} \\
\\ 
\\ 
\texttt{Wat is je volgende opdracht?} \\
            }
        }
    }
    & \\ \\

    \theutterance \stepcounter{utterance}  
    & \multicolumn{4}{p{0.6\linewidth}}{
        \cellcolor[rgb]{0.9,0.9,0.9}{
            \makecell[{{p{\linewidth}}}]{
                \texttt{\tiny{[P1$\rangle$GM]}}
                \texttt{ZEG: Klaar! Mijn C staat op (6,1), L op (4,2), P op (3,5). Alle objecten staan op de afgesproken posities!} \\
            }
        }
    }
    & & \\ \\

    \theutterance \stepcounter{utterance}  
    & & & \multicolumn{2}{p{0.3\linewidth}}{
        \cellcolor[rgb]{0.9,0.9,0.9}{
            \makecell[{{p{\linewidth}}}]{
                \texttt{\tiny{[GM$|$GM]}}
                \texttt{ZEG: Klaar! Mijn C staat op (6,1), L op (4,2), P op (3,5). Alle objecten staan op de afgesproken posities!} \\
            }
        }
    }
    & & \\ \\

    \theutterance \stepcounter{utterance}  
    & & \multicolumn{4}{p{0.6\linewidth}}{
        \cellcolor[rgb]{0.9,0.9,0.9}{
            \makecell[{{p{\linewidth}}}]{
                \texttt{\tiny{[GM$\rangle$P2]}}
                \texttt{Je bericht is doorgestuurd naar de andere speler.} \\
\\ 
\\ 
\texttt{Je speelt momenteel ronde 7 van maximaal 12.} \\
\texttt{Jullie hebben gezamenlijk 7 van 8 straffen verzameld.} \\
\texttt{De andere speler stuurde dit bericht:} \\
\\ 
\\ 
\texttt{"Klaar! Mijn C staat op (6,1), L op (4,2), P op (3,5). Alle objecten staan op de afgesproken posities!"} \\
\\ 
\\ 
\texttt{Wat is je volgende opdracht?} \\
            }
        }
    }
    & \\ \\

    \theutterance \stepcounter{utterance}  
    & & & \multicolumn{4}{p{0.6\linewidth}}{
        \cellcolor[rgb]{0.9,0.9,0.9}{
            \makecell[{{p{\linewidth}}}]{
                \texttt{\tiny{[GM$\langle$P2]}}
                \texttt{ZEG: klaar!} \\
            }
        }
    }
     \\ \\

    \theutterance \stepcounter{utterance}  
    & & & \multicolumn{2}{p{0.3\linewidth}}{
        \cellcolor[rgb]{0.9,0.9,0.9}{
            \makecell[{{p{\linewidth}}}]{
                \texttt{\tiny{[GM$|$GM]}}
                \texttt{ZEG: klaar!} \\
            }
        }
    }
    & & \\ \\

    \theutterance \stepcounter{utterance}  
    & & \multicolumn{4}{p{0.6\linewidth}}{
        \cellcolor[rgb]{0.9,0.9,0.9}{
            \makecell[{{p{\linewidth}}}]{
                \texttt{\tiny{[P1$\langle$GM]}}
                \texttt{Je bericht is doorgestuurd naar de andere speler.} \\
\\ 
\\ 
\texttt{Je speelt momenteel ronde 7 van maximaal 12.} \\
\texttt{Jullie hebben gezamenlijk 7 van 8 straffen verzameld.} \\
\texttt{De andere speler stuurde dit bericht:} \\
\\ 
\\ 
\texttt{"klaar!"} \\
\\ 
\\ 
\texttt{Wat is je volgende opdracht?} \\
            }
        }
    }
    & \\ \\

    \theutterance \stepcounter{utterance}  
    & \multicolumn{4}{p{0.6\linewidth}}{
        \cellcolor[rgb]{0.9,0.9,0.9}{
            \makecell[{{p{\linewidth}}}]{
                \texttt{\tiny{[P1$\rangle$GM]}}
                \texttt{ZEG: klaar!} \\
            }
        }
    }
    & & \\ \\

    \theutterance \stepcounter{utterance}  
    & & & \multicolumn{2}{p{0.3\linewidth}}{
        \cellcolor[rgb]{0.9,0.9,0.9}{
            \makecell[{{p{\linewidth}}}]{
                \texttt{\tiny{[GM$|$GM]}}
                \texttt{ZEG: klaar!} \\
            }
        }
    }
    & & \\ \\

    \theutterance \stepcounter{utterance}  
    & & & \multicolumn{2}{p{0.3\linewidth}}{
        \cellcolor[rgb]{0.9,0.9,0.9}{
            \makecell[{{p{\linewidth}}}]{
                \texttt{\tiny{[GM$|$GM]}}
                \texttt{true} \\
            }
        }
    }
    & & \\ \\

    \theutterance \stepcounter{utterance}  
    & & & \multicolumn{2}{p{0.3\linewidth}}{
        \cellcolor[rgb]{0.9,0.9,0.9}{
            \makecell[{{p{\linewidth}}}]{
                \texttt{\tiny{[GM$|$GM]}}
                \texttt{Initial states:} \\
\texttt{```} \\
\texttt{Player 1:} \\
\texttt{ 1234567} \\
\texttt{╔══╤════╗} \\
\texttt{║◌◌│C◌◌◌║ 1} \\
\texttt{╟──┤◌L┌─╢ 2} \\
\texttt{║◌◌├──┤◌║ 3} \\
\texttt{╟──┘◌◌├─╢ 4} \\
\texttt{║◌◌P◌◌│◌║ 5} \\
\texttt{╟─┐◌┌─┼─╢ 6} \\
\texttt{║◌│◌│◌│◌║ 7} \\
\texttt{╚═╧═╧═╧═╝} \\
\\ 
\\ 
\texttt{Player 2:} \\
\texttt{ 1234567} \\
\texttt{╔══╤════╗} \\
\texttt{║◌◌│L◌◌C║ 1} \\
\texttt{╟──┤◌◌┌─╢ 2} \\
\texttt{║◌◌├──┤◌║ 3} \\
\texttt{╟──┘◌◌├─╢ 4} \\
\texttt{║◌◌P◌◌│◌║ 5} \\
\texttt{╟─┐◌┌─┼─╢ 6} \\
\texttt{║◌│◌│◌│◌║ 7} \\
\texttt{╚═╧═╧═╧═╝} \\
\\ 
\texttt{```} \\
            }
        }
    }
    & & \\ \\

    \theutterance \stepcounter{utterance}  
    & & & \multicolumn{2}{p{0.3\linewidth}}{
        \cellcolor[rgb]{0.9,0.9,0.9}{
            \makecell[{{p{\linewidth}}}]{
                \texttt{\tiny{[GM$|$GM]}}
                \texttt{End states:} \\
\texttt{```} \\
\texttt{Player 1:} \\
\texttt{ 1234567} \\
\texttt{╔══╤════╗} \\
\texttt{║◌◌│◌◌C◌║ 1} \\
\texttt{╟──┤L◌┌─╢ 2} \\
\texttt{║◌◌├──┤◌║ 3} \\
\texttt{╟──┘◌◌├─╢ 4} \\
\texttt{║◌◌P◌◌│◌║ 5} \\
\texttt{╟─┐◌┌─┼─╢ 6} \\
\texttt{║◌│◌│◌│◌║ 7} \\
\texttt{╚═╧═╧═╧═╝} \\
\\ 
\\ 
\texttt{Player 2:} \\
\texttt{ 1234567} \\
\texttt{╔══╤════╗} \\
\texttt{║◌◌│◌◌C◌║ 1} \\
\texttt{╟──┤L◌┌─╢ 2} \\
\texttt{║◌◌├──┤◌║ 3} \\
\texttt{╟──┘◌◌├─╢ 4} \\
\texttt{║◌◌P◌◌│◌║ 5} \\
\texttt{╟─┐◌┌─┼─╢ 6} \\
\texttt{║◌│◌│◌│◌║ 7} \\
\texttt{╚═╧═╧═╧═╝} \\
\\ 
\texttt{```} \\
            }
        }
    }
    & & \\ \\

    \theutterance \stepcounter{utterance}  
    & & & \multicolumn{2}{p{0.3\linewidth}}{
        \cellcolor[rgb]{0.9,0.9,0.9}{
            \makecell[{{p{\linewidth}}}]{
                \texttt{\tiny{[GM$|$GM]}}
                \texttt{* success: True} \\
\texttt{* lose: False} \\
\texttt{* aborted: False} \\
\texttt{{-}{-}{-}{-}{-}{-}{-}} \\
\texttt{* Shifts: 1.00} \\
\texttt{* Max Shifts: 4.00} \\
\texttt{* Min Shifts: 2.00} \\
\texttt{* End Distance Sum: 0.00} \\
\texttt{* Init Distance Sum: 4.41} \\
\texttt{* Expected Distance Sum: 12.57} \\
\texttt{* Penalties: 7.00} \\
\texttt{* Max Penalties: 8.00} \\
\texttt{* Rounds: 7.00} \\
\texttt{* Max Rounds: 12.00} \\
\texttt{* Object Count: 3.00} \\
            }
        }
    }
    & & \\ \\

\end{supertabular}
}

\end{document}
